\documentclass[11pt, twocolumn, landscape, a4paper]{article}

\usepackage[margin=1cm]{geometry}
\usepackage[utf8]{inputenc}
\usepackage[spanish]{babel}
\usepackage{amssymb, amsmath, amsbsy}
\usepackage{tasks}
\usepackage{enumerate}
\usepackage{enumitem}
\usepackage{graphicx}

\usepackage[pdftex]{hyperref}
\hypersetup{pdftitle={Repaso en Polinomios / Valor numérico},pdfauthor={Luis Carrillo Gutiérrez}}

\fontfamily{cmss}\selectfont
\renewcommand{\familydefault}{\sfdefault}
\pagestyle{empty}

\begin{document}
\subsection*{Polinomios}

\begin{itemize}
\item{Dados los polinomios: \quad$P(x) = 3x + 2$; \quad$Q(x + 1) = 2x - 5$\\ Hallar el valor numérico en las siguientes expresiones:
\begin{itemize}[label=$\circ\;$]
\item{$P(x - 1) + Q(x + 2)$}
\item{$P(Q(x))$}
\item{$Q(P(x)) - 10$}
\item{$P(P(x + 1))$}
\item{$P(2 - x) - Q(-x)$}
\item{$P(4x) + Q(7 - x)$}
\item{$2P(2x) - 3Q(x - 3)$}
\item{$P(P(P(x))) - 8$}
\end{itemize}
}
\item{De un polinomio $P(x)$ de segundo grado se sabe que : $P(1) = 4$; $P(0) = 5$; $P(-2) + P(2) = 26$, entonces el valor de $P(3)$ viene a ser :
\begin{tasks}(5)
\task{10}\task{11}\task{12}
\end{tasks}
}
\item{Sean las expresiones : \quad$F(\sqrt{x} + 1) = x^{2} + x + 2$; \quad$G(\dfrac{2x - 3}{x + 4}) = 3x + 1$. \\ Calcular : $F(4) + G(1)$
\begin{tasks}(5)
\task{96}\task{104}\task{116}
\end{tasks}
}
\end{itemize}
\end{document}
